% Folder 79, batch 13 — archivists' pages 241–260.
% Pass: Opus 5.5 (claude-opus-5-5), 2026-10-03 — first pass, unchecked against the pages by a human.
%
% Notation fixed for this batch: \mathfrak{S}_S for the symmetric group of S;
% \mathfrak{P}_{\{1,2\}}(S) for the parts with one or two elements; \mathfrak{z},
% \mathfrak{z}_H for his curly « 3 » (kernel / stabiliser of the oriented edge);
% \widetilde{Z} = G/H; T_0, T_1 for his T_{0G}, T_{1G} where the index G is struck.
% Equation numbers (65)–(134) are his own; his pagination runs 124–133.
\documentclass[11pt,a4paper]{article}
\input{../preamble/grothendieck}

\folder{79}
\batch{13}
\pages{241}{260}
\foldertitle{2-polyèdres réguliers triangulés, et modules libres de rangs 2 sur les anneaux $\Lambda=\mathbb{Z}/n\mathbb{Z}$ et $\Lambda=\mathbb{Z}$ : notes manuscrites (s.d.).}
\dating{[à partir de 1980]}
\watermark{Édition de démonstration}

\begin{document}

\page{241}
\note{la phrase commence sur la page précédente, hors de ce lot ; les numéros (62), (64) renvoient à des formules antérieures.}

et toutes les situations de départ (62), en termes de \struck{\ill{}} données de théorie des
groupes plus explicites ou « élémentaires ».

Quand on part de l'ensemble $S$, ou les deux sous-groupes
\[
(65)\qquad H,\ Z \subset \mathfrak{S}_S
\]
\note{après le $H$ de (65), un signe biffé, peut-être un tilde.}
qui commutent l'un à l'autre, $Z$ commutatif, alors \struck{la} \add{une} structure de
\add{multi}graphe \add{$(S, A = \coprod A_i)$} sur $S$, admettant $H$ comme groupe
d'automorphismes transitif sur \add{l'ens. des} $\vec{A}_i$ \add{($i\in\pi$)}, et $Z$ comme groupe
d'homothéties, est connue quand on connaît \add{(65) et}
$A \;(= \bigcup_{i\in\pi} A_i) \subset \mathfrak{P}_{\{1,2\}}(S)$,
\[
(66)\qquad A \subset \underbrace{\mathfrak{P}_{\{1,2\}}(S)}_{\text{parties à 1 ou 2 élém.}} ,
\]
\struck{qui doit être nécessairement un espace homogène sous le groupe} [car on récupère
alors les $A_i$ ($i\in\pi$) comme les orbites de $H$ dans $A$, et \ill{}
\[
\text{\struck{$(67)\quad \widetilde{H} = H.Z \subset \mathfrak{S}_S$}}
\]
\[
(67')\qquad \pi \simeq A/H .
\]

\struck{Pour expliciter les conditions sur les données (65) (66) pour qu'elles correspondent
bien à un multigraphe à groupe d'homothéties homogène}, \struck{$Z$ étant groupe
d'automorphismes homogène}, \struck{$A$ on trouve que} \struck{$\lambda\in Z$} \struck{$A$
stable par} \ill{} \struck{$(S,A)$}, \struck{$\{s,t\}\in A$} \struck{$\Rightarrow \{\lambda s, t\}\in A$} \struck{\ill{}}
\note{tout ce dernier passage est biffé de traits obliques et encadré ; plusieurs
mots en restent illisibles, les fragments lus sont donnés dans l'ordre de la page.}

\page{242}
\note{p. 124 de l'auteur.}

Il serait facile d'expliciter les conditions sur les données (65) (66) pour qu'elles
correspondent à un multigraphe à homothéties et qui d'autre part soit homogène — je préfère
plutôt écrire les conditions pour que les données
\[
(68)\qquad H,\ Z \subset \mathfrak{S}_S \quad\text{et}\quad (s_0,t_0)\in S\times S
\]
correspondent à un multigraphe à homothéties et qui soit orienté homogène — en notant que
\struck{\ill{}} la donnée de $(s_0,t_0)\in\vec{A}$ \struck{\ill{}}, en plus de $H$, $Z$, permet de
récupérer $A$, comme
\[
(69)\qquad A = \bigl\{\{u(s_0),\, u(\lambda t_0)\} \bigm| \lambda\in Z,\ u\in H \bigr\} .
\]
\note{lecture de l'intérieur de l'accolade incertaine, la plume y repasse.}

On trouve le système de conditions
\[
(70)\qquad
\begin{cases}
\text{a) } H \text{ transitif sur } S,\ H \text{ et } Z \text{ commutent.}\\
\text{b) } \forall \lambda\in Z,\ \exists u\in H \text{ tel que } u(s_0,t_0) = (s_0, \lambda t_0) .
\end{cases}
\]
\note{dans b), un signe est surchargé devant $s_0$.}

On récupère alors $A_0$
\[
(71)\qquad A_0 = \bigl\{\{us_0,\, ut_0\} \bigm| u\in H\bigr\} ,
\]
et par b) $Z$ sera un groupe d'homothéties, \uncertain{ce qui assure} un multigraphe pour
$(S, A_0)$, ce qui \uncertain{garantit} \ill{} à homothéties $(S,A)$.

On n'exclut pas le cas \ill{} où on aurait $s_0=t_0$, auquel cas $A_0$ est formé de toutes les
parties ponctuelles de $S$, et la condition b) signifie que $\lambda\in Z \Rightarrow \lambda^2=1$. Si $\lambda\in Z$, $\lambda\neq 1$, alors $A_\lambda$ est formé des orbites sous
$\lambda$, qui sont des parties à deux éléments. \struck{l'étoile}

\page{243}
Si on ne fait pas d'hypothèse supplémentaire, il n'est pas sûr que l'opération de $Z$ sur
$\pi$ soit fidèle, i.e. qu'on ait
\[
(72)\qquad
\begin{cases}
\forall \lambda\in Z \text{\struck{ tel qu'il existe}}\ u\in H \text{\struck{ avec}} :\\
(s_0, \lambda t_0) = u(s_0,t_0) \;\Longrightarrow\; \lambda = 1 \\
\text{i.e. } us_0 = s_0,\ u(t_0) = \lambda t_0 .
\end{cases}
\]

Pour ceci, il suffit que tout \uncertain{sommet} soit déterminé par son étoile \ill{} condition
qui me semble, et je pense s'exprimer \ill{}, il me semble, \ill{} \uncertain{plutôt}
\uncertain{particulièrement} plus en termes des données
$S$, $H\subset\mathfrak{S}_S$, $(s_0,t_0)\in S\times S$.

Je vais, pour simplifier, supposer qu'on a (72), \uncertain{comme c'est le cas} de
\struck{\ill{}} $G$, — dans (64), l'homo \ill{} non dans $Z$, mais dans le groupe quotient
$Z/\mathfrak{z}_Z$, où $\mathfrak{z}_Z$ est le noyau de l'opération de $Z$ sur $\pi$, i.e. le
sous-groupe formé des $\lambda\in Z$ pour lesquels $\exists u\in H$ avec $u(s_0,t_0) = (s_0, \lambda t_0)$.
\note{il écrit d'une sorte de « 3 » bouclé ce groupe et le suivant ; transcrit
$\mathfrak{z}$ dans tout le lot.}

Posons
\[
(73)\qquad
\begin{cases}
U_0 = H_{s_0} = \{u\in H \mid us_0 = s_0\}\\
U_1 = H_{t_0} = \{u\in H \mid ut_0 = t_0\}\\
\mathfrak{z}_H = U_0\cap U_1 = H_{(s_0,t_0)} = \{u\in H \mid us_0 = s_0,\ ut_0 = t_0\} ,
\end{cases}
\]
et soit
\[
(74)\qquad \sigma\in H \qquad \sigma(s_0,t_0) = (t_0,s_0) \quad\text{i.e.}\quad \sigma s_0 = t_0,\
\sigma t_0 = s_0 .
\]
On a donc
\[
(75)\qquad
\begin{cases}
U_1 = \sigma(U_0),\quad U_0 = \sigma(U_1)\\
\sigma^2 \in \mathfrak{z}_H \stackrel{\text{déf}}{=} U_0\cap U_1 .
\end{cases}
\]

\page{244}
\note{p. 125 de l'auteur.}

On voit tout de suite

(76) La donnée des graphes à opérateurs homogènes $(S,\ A_0\subset\mathfrak{P}_{\{1,2\}}(S),\ H\subset\mathfrak{S}_S)$, munis de l'arête orientée $(s_0,t_0)$, équivaut à la donnée des groupes
$H$, des deux sous-groupes $U_0$, $U_1$ de $H$, et de $\sigma\in H$ satisfaisant les relations
(75), et de plus
\[
(77)\qquad \bigcap_{u\in H} \operatorname{int}(u)\, U_0 = \{1\}
\]
qui exprime que l'opération de $H$ sur l'espace homogène $H/U_0$ est fidèle.

On récupère la situation géométrique, ou la situation algébrique, par
\[
(78)\qquad
\begin{cases}
\underline{S} \simeq H/U_0, \quad s_0 = U_0,\ t_0 = \sigma.U_0 \ \text{(translaté à gauche)}\\
\vec{A}_0 \simeq H/\mathfrak{z} .
\end{cases}
\]

\struck{Dans} L'énoncé (76) n'est correct que si, \ill{} sous-entendu que pour les \ill{}
\ill{} le choix de $\sigma$ — en termes géométriques il y a aussi le choix de $\sigma$
satisfaisant (75), \uncertain{soit} que dans la \uncertain{structure} \add{(qui implique
$\mathfrak{z}_H$ uniquement)} \struck{des données} algébrique, $\sigma$ \ill{} « modulo
$\mathfrak{z}_H$ » (i.e. que ce n'est pas $\sigma$ \ill{} qui est donné, c'est la classe
$\sigma\mathfrak{z}_H = \mathfrak{z}_H\sigma$ dans $\operatorname{Norm}_H(\mathfrak{z}_H)/ \mathfrak{z}_H$). Cette \uncertain{ambiguïté} disparaît quand $H$ est \emph{simplement}
transitif sur l'ens. $\vec{A}_0$ des arêtes orientées : \struck{\ill{}}
\[
(79)\qquad H \text{ simplement transitif sur l'ens. } \vec{A}
\;\Longleftrightarrow\; \mathfrak{z}_H \;(\stackrel{\text{déf}}{=} U_0\cap U_1) = \{1\} ,
\]
\struck{\ill{}} condition qui implique l'unicité de $\sigma$, et
\[
(80)\qquad \sigma^2 = 1 .
\]

\page{245}
Nous allons expliciter maintenant la donnée de $Z$, en introduisant le sous-groupe de $H$
\[
(80)\qquad T_H \stackrel{\text{déf}}{=} \{u\in H \mid \exists\lambda\in Z,\ us_0 = \lambda s_0,\
ut_0 = \lambda^{-1}t_0\} .
\]
\note{le numéro (80) est employé deux fois, p. 244 et ici.}

C'est bien un sous-groupe, \struck{\ill{}} si $u$ correspond à $\lambda$, $v$ correspond à
$\mu$, alors $uv$ correspond à $\mu\lambda$. \struck{D'ailleurs le $\lambda$ associé à $u$ est
unique, car $Z$ opère librement sur $S$} \struck{\ill{}} \add{mieux} \struck{vu que}
\[
\text{\struck{$(\lambda s_0, \lambda^{-1}t_0) = (s_0,t_0) \Longrightarrow \lambda = 1$}}
\]
\struck{en \ill{} de (72)}. On a donc un homomorphisme de groupes canonique
$T_H\to Z$, dont le noyau est $\mathfrak{z}$ \add{et \ill{} l'hom. $\delta\vert Z$}
\[
(81)\qquad 1 \longrightarrow \mathfrak{z}_H \longrightarrow T_H \xrightarrow{\;k\;} Z
\longrightarrow 1
\]
qui dans le cas simplement transitif s'écrit
\[
(81')\qquad T_H \xrightarrow{\;\sim\;} Z .
\]

Notons
\[
(82)\qquad T_H \text{ normalise } U_0 \text{ et } U_1, \qquad
T_H\cap U_0 = T_H\cap U_1 = \mathfrak{z}_H
\]
\note{« et $U_1$ » est ajouté au-dessus de la ligne.}
\note{suivent deux lignes hachurées, illisibles.}

En effet, si $t\in T_H$, $u\in U_0$, on a
\[
tut^{-1}(s_0) = t\,\underbrace{u\lambda^{-1}(s_0)}_{\lambda^{-1}\underbrace{u(s_0)}_{s_0}}
= t(\lambda^{-1}s_0) = \lambda^{-1}ts_0 = \lambda^{-1}\lambda s_0 = s_0
\]
d'où $tut^{-1}\in U_0$, \struck{\ill{}} \add{par symétrie on voit} que $T_H$ normalise $U_1$.
Si $t\in T_H\cap U_0$, on a
\[
t(s_0,t_0) = (\lambda s_0, \lambda^{-1}t_0) = (s_0, \lambda^{-1}t_0)
\]
puisque \uncertain{$ts_0 = s_0$},

\page{246}
\note{p. 126 de l'auteur.}

d'où $\lambda = 1$ par (72), d'où $t(s_0,t_0) = (s_0,t_0)$ i.e. $t\in\mathfrak{z}_H$, donc
$T_H\cap U_0\subset\mathfrak{z}_H$, l'inclusion inverse étant triviale. \struck{Notons} Posons
alors
\[
(83)\qquad
\begin{cases}
B_{0H} = T_H.U_0\\
B_{1H} = T_H.U_1
\end{cases}
\qquad \text{(produits ½-directs si } \mathfrak{z}_H = \{1\}) ,
\]
je dis qu'on a
\[
(84)\qquad
\begin{cases}
B_{0H} = \{u\in H \mid us_0\in Z.s_0\} = H_{D_0}\\
B_{1H} = \{u\in H \mid us_1\in Z.t_0\} = H_{D_1} ,
\end{cases}
\]
\note{« $us_1$ » sic, pour $ut_0$.}
où $D_0$, $D_1$ sont les « voisinages » de $s_0$, $t_0$, i.e. leurs orbites sous $Z$. La
vérification est immédiate, sur le fait que l'hom $k : T_H\to Z$ est épi.

Notons aussi
\[
(85)\qquad B_{0H}\cap B_{1H} = T_H, \qquad \text{d'où aussi}\quad
B_{0H}\cap U_1 = U_0\cap B_{1H} = \mathfrak{z}_H .
\]
En effet si $t\in B_{0H}\cap B_{1H}$, i.e. si $t\in H$ tel que l'on ait
\[
t(s_0,t_0) = (\lambda s_0, \mu t_0) \qquad \lambda,\mu\in Z ,
\]
$(\lambda s_0,\mu t_0)$ est une $\lambda\mu$-arête, et $t\in H$, comme \ill{} est une
$A_0$-arête, on a $\lambda\mu = 1$, i.e. $\mu = \lambda^{-1}$, d'où $t\in T$. \struck{\ill{}}

Notons maintenant que les groupes $U_0$, $U_1$, $T_H$, $B_{0H}$, $B_{1H}$ dépendent du choix
de $(s_0,t_0)\in\vec{A}_0$ ; mais $U_0$ et $U_1$ s'interchangent quand on échange $s_0$ et
$t_0$, ainsi que $B_{0H}$ et $B_{1H}$, tandis que $T_H$ reste inchangé, il ne dépend que de
\emph{l'arête} $\{s_0,t_0\}$ — mais $k$ est remplacé par $k' = t\mapsto k(t)^{-1}$.

\page{247}
En termes de $\sigma\in H$, cela s'exprime par
\[
(86)\qquad
\begin{cases}
\sigma \text{ normalise } T_H, \text{ et } \forall t\in T_H, \text{ on a}\\
\qquad \sigma(t) \equiv t^{-1} \ (\mathfrak{z}_H)\\
\text{donc}\quad \sigma(t) = t^{-1} \ \text{si}\ \mathfrak{z}_H = \{1\} .
\end{cases}
\]
\note{avant « $\forall t$ », quelques mots biffés, illisibles.}
\[
(87)\qquad B_{1H} = \sigma(B_{0H}), \qquad B_{0H} = \sigma(B_{1H}) .
\]

Notons maintenant

(88) La donnée, \add{*} pour un graphe homogène défini par $(H,\ U_0\subset H,\ \sigma\in H)$
(satisfaisant (75) : $\sigma^2\in U_0\cap\sigma(U_0) = \mathfrak{z}$) \add{et satisfaisant à
(72)*}, d'un groupe d'homothéties $Z$, équivaut à celle d'un sous-groupe $T_H\subset H$,
normalisant $U_0$, tel qu'on ait $T_H\cap U_0 = \mathfrak{z}_H$, [que $T_H/\mathfrak{z}_H$ soit
commutatif], \struck{et} que $\sigma$ \emph{normalise} $T_H$, $\sigma(t) = t^{-1}$
\add{mod $\mathfrak{z}_H$} \struck{\ill{}}, pour $t\in T_H$, et \struck{tant} satisfasse et tel
enfin qu'on ait $[T U_0]\cap U_1 = \mathfrak{z}_H$ (d'où \struck{\ill{}} $T_H\cap U_0 = \mathfrak{z}_H$).

\marginal{* i.e. une opération \emph{fidèle} de $Z$ sur l'ens. $\pi$ des structures de graphe
engendrées par $(S, \vec{A}_0)$, sous $Z$}

\marginal{NB On voit alors que [ ] résulte des autres}

En effet, si $T_H$ \add{$B_{0H}$} \ill{} \add{$\stackrel{\text{déf}}{=} Z$} sous-groupe de $H$
normalisant $U_0$, alors
\[
T_H / U_0\cap T_H \simeq \underbrace{T_H.U_0}_{B_{0H}} / U_0
\]
opère \add{fidèlement} sur l'espace homogène $S\simeq H/U_0$ \add{$B_{0H}$}, il faut vérifier
que c'est un groupe d'homothéties, \uncertain{soit} \ill{} commutativité, sachant \ill{}
pour $t\in T_N$ on a
\[
\sigma(t).t \in \mathfrak{z}_H \subset U_1\cap U_0 \subset U_1.U_0 ,
\]
\ill{} déjà fait et immédiat aussi aussitôt que \ill{} elle implique que $t$ est une
homothétie. \struck{À} \struck{Le cas \ill{}, on a utilisé seulement l'inclusion} \struck{$T_H\cap U_0\subset\mathfrak{z}$}

\page{248}
\note{p. 127 de l'auteur.}

La condition $[TU_0]\cap U_1$ ne fait alors que exprimer (72)*.

La situation est particulièrement jolie dans le cas où $\mathfrak{z}_H = 1$ :

(89) Les données, pour un graphe homogène, \emph{simplement} \emph{transitif} $H$ sur
$\vec{A}_0$, \ill{} groupe défini par $(H,\ U_0\subset H,\ \sigma\in H)$, d'un groupe
d'homothéties $Z$ satisfaisant la condition de fidélité (72), équivaut à celle d'un sous-groupe
$T_H$ de $H$ satisfaisant les conditions
\begin{itemize}
\item \struck{$T_H$ normalise $U_0$, \ill{}} \struck{$T_H\cap U_0 = \{1\}$}
\item a) $\sigma$ normalise $T_H$ et y induit $t\mapsto t^{-1}$ \add{(donc aussi $U_1 = \sigma(U_0)$)}
\item b) \struck{\ill{}} $T_H$ normalise $U_0$
\item c) l'application
\[
U_0\times T_H\times U_1 \longrightarrow H, \qquad (u_0,t,u_1)\longmapsto u_0tu_1
\]
est \emph{injective}.
\end{itemize}
\note{dans la marge droite, à hauteur de a), un bloc de formules raturé, illisible.}

Je vais maintenant expliciter, pour une telle donnée, la construction des diagrammes (64) — je
me bornerai au cas $H$ simplement transitif sur $\vec{A}_0$, pour simplifier, cas tellement plus
joli que le cas général !

J'introduis les sous-groupes
\[
\text{\struck{$(90)\quad H_{?} = G_{s_0} = \{g\in G \mid \ldots\}$}}
\]
\note{dans cette ligne biffée, l'indice de $H$ et la condition sont illisibles.}
\[
(90)\qquad T_G = G_{(D_0,D_1)} = \{g\in G \mid \exists\lambda,\mu\in Z,\ gs_0 = \lambda s_0,\
gt_0 = \mu t_0\}
\]
\note{sous la condition, biffé : « i.e. $gD_0 = D_0$, $gD_1 = D_1$ ».}
Les éléments $\lambda,\mu\in Z$ sont bien déterminés \add{par $g$}, d'où un homomorphisme
$T_G\to Z\times Z$, \struck{\ill{}} et on trouve un homomorphisme

\page{249}
évidemment surjectif, dont le noyau est le groupe
\[
(91)\qquad \mathfrak{z} = G_{(s_0,t_0)} = \{g\in G \mid g(s_0) = s_0,\ g(t_0) = t_0\}
\]
\marginal{on a $\mathfrak{z}\subset$ \uncertain{$G_0(s_0,t_0)$}}
d'où une suite exacte
\[
(92)\qquad
\begin{gathered}
1 \longrightarrow \mathfrak{z} \longrightarrow T_G \longrightarrow Z\times Z \longrightarrow 1\\
g \longmapsto (\delta_0(g), \delta_1(g))
\end{gathered}
\]
On pose
\[
(93)\qquad
\begin{cases}
T_0 = \operatorname{Ker}\delta_0 = \{g\in G \mid gs_0 = s_0,\ gt_0\in Zt_0\}\\
T_1 = \operatorname{Ker}\delta_1 = \{g\in G \mid gs_0\in Zs_0,\ gt_0 = t_0\}
\end{cases}
\]
\note{il écrit d'abord $T_{0G}$, $T_{1G}$, puis biffe l'indice $G$ partout sur cette page et
les suivantes ; on transcrit $T_0$, $T_1$. Dans (93), $\operatorname{Ker}\delta_0$ est suivi
d'un terme biffé illisible. Les définitions de $T_0$, $T_1$ semblent échangées par rapport à
$\operatorname{Ker}\delta_0$, $\operatorname{Ker}\delta_1$ : transcrit tel quel.}

Notons que
\[
(94)\qquad
\begin{cases}
T_0,\ T_1 \subset T_G, \qquad Z\subset T_G ,\\
T_G = T_0.T_1 = T_0.Z = Z.T_1\\
T_0\cap T_1 = T_0\cap Z = Z\cap T_1 = \mathfrak{z}
\end{cases}
\]
et
\[
(95)\qquad
\begin{cases}
\sigma(T_G) = T_G\\
\sigma(T_0) = T_1, \quad \sigma(T_1) = T_0\\
\text{et bien sûr}\quad \sigma(Z) = Z \quad (Z \text{ central ?}) .
\end{cases}
\]

On a aussi
\[
(96)\qquad
\begin{gathered}
T_G\cap H = T_H \subset \operatorname{Ker}(\delta_0.\delta_1 : T_G\to Z)\\
t \longmapsto \delta_0(t)\delta_1(t)
\end{gathered}
\]
\struck{et \ill{} $\delta_H : T_H \xrightarrow{\sim} Z$ est \ill{} 2 : 1} \struck{Introduisons}
\struck{\ill{} $T_G$ normalise $U_0$ et centralise $T_H$}
\[
\text{\struck{$(97)\quad \operatorname{Ker}(\delta_0\delta_1 : T_G\to Z) = T_H$}}
\]
\struck{(97) Comme $T_G$ est défini comme} \struck{\ill{} le centralisateur}
\note{passage barré de plusieurs traits et d'une ligne ondulée.}

La formule (96) implique que $T_G$ normalise $T_H$ (puisque $H$ est distingué dans $G$), donc
comme $T_G/\mathfrak{z}$ est commutatif et $T_H\cap\mathfrak{z}\subset H\cap\mathfrak{z} = \{1\}$, on voit
\[
(97)\qquad T_G \text{ centralise } T_H
\]

\page{250}
\note{p. 128 de l'auteur.}

(\uncertain{cela} qu'il n'est pas sûr que $T_G$ soit commutatif, car il n'est pas sûr que
$\mathfrak{z}$ le soit…). On voit aussi immédiatement
\[
(98)\qquad
\begin{cases}
T_G \text{ normalise } U_0 \text{ et } U_1, \quad T_G\cap H = T_H \ \text{d'où}\\
T_G\cap U_0 = T_G\cap U_1 = \{1\} \quad (\text{car } T_H\cap U_0 = T_H\cap U_1 = \{1\})
\end{cases}
\]
On introduit donc les sous-groupes
\[
(99)\qquad
\left.
\begin{aligned}
B_{0G} &= T_G.U_0\\
B_{1G} &= T_G.U_1
\end{aligned}
\right\}\ \text{produits semi-directs}
\]
\note{dans (99), un petit trait au-dessus de chaque $B$, peut-être un prime ; il disparaît
dans la suite.}
donc
\[
(100)\qquad \sigma(B_{0G}) = B_{1G}, \qquad \sigma(B_{1G}) = B_{0G} .
\]

On prouve
\[
(101)\qquad
\begin{cases}
B_{0G} = G_{D_0} = \{g\in G \mid gs_0\in Zs_0\}\\
B_{1G} = G_{D_1} = \{g\in G \mid gt_0\in Zt_0\}
\end{cases}
\]
d'où
\[
(102)\qquad B_{0G}\cap B_{1G} = T_G .
\]
Pour prouver (101), il suffit de prouver que si $g\in G$, \struck{\ill{}} $\lambda\in Z$ sont
tels que
\[
gs_0 = \lambda s_0
\]
alors $g\in T_GU_0$. Or quitte à corriger \add{$g$} par un élément convenable de $T_G$, OPS
$\delta(g) = 1$ i.e. $g\in G_0$. \struck{Quitte à} quitte à corriger par un élément de
$\mathfrak{z}\subset T_G$, OPS $g\in H$ (cf plus bas \ill{}). Mais on a alors par (84)
$g\in B_{0H} \stackrel{\text{déf}}{=} T_H.U_0 \subset T_G.U_0$ cqfd.

On a utilisé en passant
\[
(103)\qquad G_0 = \mathfrak{z}.H \qquad (\text{½ direct}),
\]

\page{251}
qui résulte aussitôt des définitions, et du fait que $H$ est simpl. transitif sur $\vec{A}_0$.
Notons aussi
\[
(104)\qquad
\begin{gathered}
U_0\times T_G\times U_1 \longrightarrow G \qquad \text{est injectif}\\
(u_0,t,u_1) \longmapsto u_0tu_1
\end{gathered}
\]
en effet, comme $T_G$ normalise $U_0$, $U_1$, cela revient aux formules
\[
T_G\cap U_0 = T_G\cap U_1 = \{1\} \quad (98) \qquad \text{et}
\]
\[
(*)\qquad \underbrace{U_0.T_G}_{B_{0G}} \cap\, U_1 = \{1\}
\]
Mais on a
\[
(105)\qquad
\begin{cases}
B_{0G}\cap H = B_{0H} \quad (= U_0.T_H)\\
B_{1G}\cap H = B_{1H} \quad (= U_1.T_H)
\end{cases}
\]
(résulte immédiatement de $T_G\cap H = T_H$), donc $(*)$ se réduit \struck{\ill{}} encore à : la
formule analogue dans $H$
\[
U_0.T_H\cap U_1 = \{1\} ,
\]
qui est satisfaite par hypothèse.

On a un peu bousculé, sans doute : la construction de (64) en termes des données
$(H, U_0, \sigma, T_H)$ — je vais y revenir, en introduisant
\[
(106)\qquad
\begin{cases}
B'_0 = G_{s_0} = \{g\in G \mid gs_0 = s_0\} \subset B_{0G} = G_{D_0}\\
B'_1 = G_{t_0} = \{g\in G \mid gt_0 = t_0\} \subset B_{1G} = G_{D_1} .
\end{cases}
\]
On trouve en fait, plus précisément
\[
(107)\qquad
\begin{cases}
B'_0 = T_0.U_0\\
B'_1 = T_1.U_1
\end{cases}
\qquad \text{produits semi-directs}
\]
\note{ici aussi l'indice $G$ de $B'_{0G}$, $T_{0G}$ est biffé.}

\page{252}
\note{p. 129 de l'auteur.}

et de même
\[
(108)\qquad G = T_0.H = T'_0.H \qquad (\text{produits ½ directs}) ,
\]
\note{« $T'_0$ » : lecture du second facteur incertaine, on attendrait $T_1$.}
pour prouver (108), notons d'abord
\[
T_0\cap H = \text{\struck{$T_G\cap H$}}\ \ldots = T_H
\]
\[
T_0\cap H = T_0\cap T_H = \operatorname{Ker}\bigl(T_H \xrightarrow{\;\kappa\;} T_G
\xrightarrow{\;\delta_0\;} Z\bigr) = \{1\} ,
\]
\note{au-dessus de la flèche $\kappa$, le renvoi « (96) » ; sous le diagramme, une flèche
courbe étiquetée $\kappa$ va de $T_H$ à $Z$. La première ligne est en partie biffée.}
on voit de même que $T'_0\cap H = \{1\}$. Pour voir que tout $g\in G$ s'écrit sous la forme
$t_0.u$ ($t_0\in T_0$, $u\in H$), ou ce qui revient au même, sous la forme $u.t_0$ (puisque
$T_0$ normalise $H$), OPS, comme $H$ est transitif sur $S$, quitte à \struck{\ill{}} corriger
$g$ par un $u\in H$, que l'on a
\[
gs_0 = s_0 ,
\]
\ill{} la formule (107). Quitte à corriger par un élément de $T_0$ (puisque
$T_0 \xrightarrow{\delta\vert T_0} Z$ surjectif) OPS $\delta(g) = 1$ i.e. $g\in G_0\cap B'_0$ et
\struck{\ill{}} OPS quitte à corriger par un $z\in\mathfrak{z}\subset T_0$, OPS
$g\in H\cap B'_0$ \struck{\ill{}}, et il suffit de \add{noter} \struck{voir} que
\[
(109)\qquad
\begin{aligned}
H\cap B'_0 &= U_0 ,\\
H\cap B'_1 &= U_1 ,
\end{aligned}
\]
ce qui achève de prouver (107) (108). On prouve de même
\[
(110)\qquad
\begin{cases}
B_{0G} = T_0.B_{0H}\\
B_{1G} = T_1.B_{1H} ,
\end{cases}
\]
notons aussi
\[
(111)\qquad
\begin{cases}
B'_0\cap G_0 = \mathfrak{z}.U_0\\
B'_1\cap G_1 = \mathfrak{z}.U_1
\end{cases}
\qquad \text{prod. ½ directs}
\]
\note{« $G_1$ » sic dans la seconde ligne.}

\page{253}
et
\[
(112)\qquad
\begin{cases}
B_{0G}\cap G_0 = \mathfrak{z}.B_{0H}\\
B_{1G}\cap G_0 = \mathfrak{z}.B_{1H}
\end{cases}
\qquad \text{produits semi-directs,}
\]
formules qui proviennent de la conjonction de (108), et du fait que par définition
\[
G_0 = \operatorname{Ker}(G \xrightarrow{\;\delta\;} Z)
\]
et que enfin $\delta$ s'explicite comme le composé
\[
G = T_i.H \xrightarrow{\;\text{can}\;} T_i \xrightarrow{\;\delta\vert T_i\;} Z
\qquad i\in\{0,1\}
\]
et que \struck{enfin} $\mathfrak{z}$ est le noyau de $\delta\vert T_i$.

Une partie des relations explicitées dans les formules précédentes sont résumées dans le
diagramme d'inclusions

\note{(113). Les inclusions sont étiquetées par le groupe quotient correspondant.}

\begin{tikzcd}[column sep=small, row sep=small, nodes={font=\scriptsize}]
U_0 \arrow[r, hook, "Z"] \arrow[d, hook, "\mathfrak{z}"'] & B_{0H} = T_H.U_0 \arrow[r, hook] \arrow[d, hook, "\mathfrak{z}"'] & H \arrow[d, hook] \\
\widetilde{U}_0 = \mathfrak{z}.U_0 \arrow[r, hook, "Z"] \arrow[d, hook, "\widetilde{Z}"'] & \widetilde{B}_0 = (\mathfrak{z}\times T_H).U_0 \arrow[r, hook] \arrow[d, hook, "\widetilde{Z}"'] & G_0 = \mathfrak{z}.H \arrow[d, hook, "\widetilde{Z}"] \\
B'_0 = T_0.U_0 \arrow[r, hook, "Z"] & B_{0G} = (T_0\times T_H).U_0 \arrow[r, hook] & G = T_0.H
\end{tikzcd}

\note{sur la page, les nœuds du milieu et du bas portent aussi les écritures
$\widetilde{B}_0 = \mathfrak{z}.B_{0H}$ et $B_{0G} = T_0.B_{0H}$.}

dont les quatre carrés \uncertain{commutatifs} sont \uncertain{cartésiens et cocartésiens}, et
où on a écrit, au-dessus des inclusions distinguées, la valeur du groupe quotient
correspondant, où j'ai posé
\[
\widetilde{Z} \stackrel{\text{déf}}{=} G/H \xrightarrow[\text{surj.}]{} G/G_0
\overset{\delta}{\hookrightarrow} Z ,
\]
composé \ill{} de noyau $\mathfrak{z}$ ; d'où un hom. can.
\[
(115)\qquad 0 \longrightarrow \mathfrak{z} \longrightarrow \widetilde{Z} \xrightarrow{\;c\;} Z
\]
\marginal{Je dis que $G\to G/H \simeq \widetilde{Z}$, quotient \ill{} (114) l'hom. can.
\ill{} d'où (115)}
\note{le numéro (114) n'apparaît que dans cette note marginale, lue de biais.}

\page{254}
\note{p. 130 de l'auteur.}

qui est surjectif ssi on est dans la « situation d'isomorphie » i.e. $\delta : G\to Z$
surjectif.

On a un diagramme symétrique (113') en y remplaçant l'indice $0$ par $1$, il se déduit de (113)
par \struck{le transformé} \add{la transformation par} l'automorphisme $\operatorname{int}(\sigma)$,
où $\sigma$ est considéré comme élément de $G$.

\marginal{NB On a aussi des iso (115') $T_0 \xrightarrow{\sim} \widetilde{Z}$, $T_1 \xrightarrow{\sim} \widetilde{Z}$ \add{(induits par la projection de $G$)}, mais \ill{}
d'identification de $T_0$, $T_1$ \ill{} $G$}

\struck{Ce n'est qu'une description « symétrique » de} \struck{$G$, on peut écrire}
\[
(116)\qquad G = T_G.H, \qquad \text{où}\quad T_G\cap H = T_H \quad (98)
\]
dans cette description \add{de $G$}, il s'agit donc de bien expliciter

(116') a) le groupe $T_G$, b) le plongement de $T_H$ dans $T_G$, et c) l'action de $T_G$ sur
$H$, \struck{enfin} ce qui permet de reconstituer l'inclusion $H\subset G$ ; si on veut de plus
reconstituer (64), il faut de plus d) expliciter l'hom d'inclusion $Z\subset T_G$, et enfin
e) expliciter l'hom $\delta\vert T_G : T_G\to Z$, pour reconstituer le diagramme (64).

Nous allons expliciter $T_G$ \struck{à l'aide} \struck{de la suite exacte} comme extension
centrale
\[
\text{\struck{$(117)\quad 1\to\mathfrak{z}\to T_0\times T_1\to T_G\to 1$}}
\]
\struck{$T_0\simeq\widetilde{Z}\times\widetilde{Z}$} \struck{on suppose pour simplifier que}
\struck{$\widetilde{T}_0$ est commutatif} \struck{$T_G\simeq T_H\times T_0\simeq Z\times\widetilde{Z}$}
\note{le bas de la page est un brouillon barré de grands traits obliques, en partie illisible.}

\page{255}
\[
(117)\qquad
1 \longrightarrow \underset{\simeq Z}{T_H}
\longrightarrow T_G \xrightarrow{\;\widetilde{\delta}\vert T_G\;} \widetilde{Z}
\longrightarrow 1
\]
\note{sous $T_H$, la page écrit $\simeq Z$ verticalement ; transcrit approximativement.}
qui est scindée par $T_0$ et $T_1$, qui donnent deux \struck{\ill{}} hom
\[
(118)\qquad \widetilde{Z} \underset{u_1}{\overset{u_0}{\rightrightarrows}} T_G
\]
\note{au-dessus de $u_0$, quelques signes biffés.}
satisfaisant la relation
\[
(119)\qquad
\underbrace{u_1(\widetilde{\lambda})}_{\substack{s_0\mapsto\lambda s_0\\ t_0\mapsto t_0}}
=
\underbrace{u_0(\widetilde{\lambda})}_{\substack{s_0\mapsto s_0\\ t_0\mapsto\lambda t_0}}
\;
\underbrace{u_{c(\widetilde{\lambda})}}_{\substack{s_0\mapsto\lambda s_0\\
t_0\mapsto\lambda^{-1}t_0}}
\]
\note{au-dessus de $u_{c(\widetilde{\lambda})}$ : « él. central de $T_G$ » ; entre les deux
membres de droite, un terme hachuré.}
qui s'écrit aussi
\[
u_1(\widetilde{\lambda})\,u_0(\widetilde{\lambda})^{-1} = u_{c(\widetilde{\lambda})}
\quad\text{ou}\quad u_0(\widetilde{\lambda})^{-1}u_1(\widetilde{\lambda})
\]
et comme les deux membres sont dans $H$, il suffit de vérifier que les valeurs dans $s_0$,
$t_0$ sont les mêmes, ce qui est immédiat.

L'un et l'autre scindage de (117) donnent deux isomorphismes
\[
(120)\qquad T_G \overset{\sim}{\rightrightarrows} Z\times\widetilde{Z} ,
\]
le passage de l'un à l'autre se faisant via l'hom $\widetilde{Z} \xrightarrow{\;k\;} Z$.
\add{donc $T_G$ commutatif ;}

Quand $\widetilde{Z}$ est commutatif, \add{cette description} se simplifie, en écrivant la suite
exacte
\[
(121)\qquad
\begin{gathered}
0 \longrightarrow \mathfrak{z} \longrightarrow
\underset{\simeq\widetilde{Z}\times\widetilde{Z}}{T_0\times T_1}
\longrightarrow T_G \longrightarrow 1\\
(t,t') \longmapsto t.t' \\
z \longmapsto (z, z^{-1})
\end{gathered}
\]
qui explicite $T_G$ comme un quotient de $T_0\times T_1 \simeq \widetilde{Z}\times\widetilde{Z}$

\page{256}
\note{p. 131 de l'auteur.}

\struck{Dans $\widetilde{Z}\times\widetilde{Z}$ on a le sous-groupe diagonal,} \struck{contenant
$\mathfrak{z}$, et le quotient \ill{}} \struck{\ill{} sous-groupe $Z\subset T_G$,} \struck{\ill{}
en d'autres termes \ill{}}
\note{quatre lignes barrées.}

La suite exacte (121) se récrit en termes de $\widetilde{T}$, et inscrit la suite (115) dans le
diagramme

\note{(122)}

\begin{tikzcd}[column sep=small, row sep=small, nodes={font=\scriptsize}]
0 \arrow[r] & \mathfrak{z} \arrow[r] \arrow[d, no head, "\parallel" description] & \widetilde{Z} \arrow[r, "c"] \arrow[d, "{\widetilde{\lambda}\mapsto(\widetilde{\lambda},\widetilde{\lambda}^{-1})}"] & Z \arrow[d, hook, "\text{incl.}"] & \\
0 \arrow[r] & \mathfrak{z} \arrow[r, "{z\mapsto(z,z^{-1})}"'] & \widetilde{Z}\times\widetilde{Z} \arrow[r, "{(\widetilde{\lambda},\widetilde{\mu})\mapsto u_0(\widetilde{\lambda})u_1(\widetilde{\mu})}"'] & T_G \arrow[r] & 1
\end{tikzcd}

\note{sous le diagramme : « NB si $z\in\mathfrak{z}$, on a $u_0(z) = u_1(z)$ ».}

qui permet de retrouver l'inclusion canonique $Z\subset T_G$, quand $c : \widetilde{Z}\to Z$
est surjectif.

Ainsi, le groupe $T_G$ et l'inclusion $T_H\subset T_G$ s'explicitent (les données a) b) de
(116')) en termes du groupe $\widetilde{Z}$ et \add{l'hom $c : \widetilde{Z}\to Z$ du}
diagramme (115), qui fournit aussi l'hom. $\delta : T_G\to Z$ (donnée e) de (116')), soit
comme le composé de $\widetilde{\delta}\vert T_G : T_G\to\widetilde{Z}$ et de $c : \widetilde{Z}\to Z$, dans la \struck{\ill{}} description (117), soit dans la description
(122)), par commutativité \add{du diagramme}

\page{257}
\note{(123)}

\begin{tikzcd}[column sep=small, row sep=small, nodes={font=\scriptsize}]
\widetilde{Z}\times\widetilde{Z} \arrow[r] \arrow[d, "{(\widetilde{\lambda},\widetilde{\mu})\mapsto\widetilde{\lambda}+\widetilde{\mu}}"'] & T_G \arrow[r] \arrow[d, "\delta"] & 1 \\
\widetilde{Z} \arrow[r, "c"] & Z . &
\end{tikzcd}

\note{l'étiquette de la flèche verticale gauche se lit $\widetilde{\lambda}+\widetilde{\mu}$
(notation additive).}

Reste à décrire, \struck{\ill{}} \uncertain{pour} l'action \ill{} (115) bien connu
\add{(donnée c))}, et l'action de $T_G$ sur $H$, i.e. de $T_0\simeq\widetilde{Z}$ et de
$T_1\simeq\widetilde{Z}$ sur $H$, l'inclusion de $Z$ dans $T_G$ (donnée d) de (116')).
\add{Cette} \struck{Je dis} \struck{inclusion} peut se décrire : à l'aide de l'hom
\[
(124)\qquad c' : Z \longrightarrow \widetilde{Z} ,
\]
\note{sous (124), un petit diagramme biffé.}
défini comme composé
\[
(125)\qquad Z \xrightarrow{\;\text{incl}\;} T_G \xrightarrow{\;\widetilde{\delta}\;}
\widetilde{Z} \quad\text{ou}\quad Z \xrightarrow{\;\text{incl}\;} T_G
\xrightarrow{\;\widetilde{\delta}\vert T_G\;} \widetilde{Z} ,
\]
qui satisfait la relation
\[
(126)\qquad
\underbrace{Z \xrightarrow{\;c'\;} \widetilde{Z} \xrightarrow{\;c\;} Z}_{\lambda\mapsto\lambda^2}
\qquad \text{i.e.}\quad c(c'(\lambda)) = \lambda^2 \quad \forall\lambda\in Z .
\]
Ainsi, quand on écrit
\[
(127)\qquad T_G = T_H\times T_0
\]
la donnée de l'inclusion $Z\subset T_G$ revient à celle d'un hom
\[
(127)\qquad Z \hookrightarrow \underbrace{T_H}_{Z}\times\underbrace{T_0}_{\widetilde{Z}} ,
\]
\note{les deux numéros sont surchargés et peu lisibles ; on lit deux fois (127).}
donné par la compatibilité

\page{258}
\note{p. 132 de l'auteur.}

\note{(128)}

\begin{tikzcd}[column sep=small, row sep=small, nodes={font=\scriptsize}]
Z \arrow[r, hook] \arrow[dr, "{\lambda\mapsto(\lambda,c'(\lambda))}"'] & T_H\times T_0 \simeq T_H\times T_0 \\
 & Z\times\widetilde{Z} \arrow[u, "{(\lambda,\widetilde{\mu})\mapsto u(\lambda)u_0(\widetilde{\mu})}"', "\wr"]
\end{tikzcd}

\[
\text{i.e.}\qquad \boxed{\lambda = u(\lambda)\,u_0(c'(\lambda))}
\]

Pour vérifier la commutativité de ce diagramme, il suffit de \struck{vérifier que} noter que
les \struck{\ill{}} deux flèches \struck{\ill{}} $Z\to T_H\times T_0$ dans (128), composées avec
\[
\underbrace{T_H\times T_0}_{Z\times\widetilde{Z}} \xrightarrow{\;\widetilde{\delta}\;}
\widetilde{Z} ,
\]
donnent le même \ill{}, les deux images d'un $\lambda$ sont congrues mod $T_H$, il suffit pour
voir qu'elles sont égales, il suffit de voir que les \struck{compositions} actions sur $s_0$,
$t_0$ sont égales i.e. égales à : $\lambda s_0$, $\lambda t_0$, ce qui est immédiat. Si on avait
utilisé l'iso $T_G\simeq T_H\times T_1$, et l'iso $Z\times\widetilde{Z}$ donné par
\[
(\lambda,\widetilde{\mu}) \longmapsto u(\lambda).u_1(\widetilde{\mu})
\quad \text{\struck{$= u(\lambda).u_0(\widetilde{\mu})\,u(c(\widetilde{\lambda}))$}}
\]
on trouve
\[
\begin{aligned}
\lambda &\underset{(128)}{=} u(\lambda).u_0(c'(\lambda))
\underset{(119)}{=} u(\lambda)\,u_1(c'(\lambda))\,
u\bigl(\overbrace{c(c'(\lambda))}^{\lambda^2\ (126)}\bigr)^{-1}\\
&= u(\lambda\lambda^{-2})\,u_1(c'(\lambda))
\end{aligned}
\]
d'où
\[
(129)\qquad \lambda = u(\lambda^{-1})\,u_1(c'(\lambda)) \qquad (\lambda\in Z)
\]
\note{au-dessus de (129), une première écriture de la même ligne, hachurée.}

Dans le cas où $\widetilde{Z}$ est commutatif et où $T_G$ est décrit par (121) i.e. la
deuxième ligne de (122), on aimerait pouvoir interpréter

\page{259}
le sous-groupe $Z$ de $T_G$ comme une sorte de sous-groupe diagonal, avec un diagramme
commutatif

\note{(130)}

\begin{tikzcd}[column sep=small, row sep=small, nodes={font=\scriptsize}]
1 \arrow[r] & \mathfrak{z} \arrow[r] \arrow[d, no head, "\parallel" description] & \widetilde{Z} \arrow[r, "c"] \arrow[d, "\text{diag}"] & Z \arrow[r] \arrow[d, hook, "\text{incl.}"] & 1 \\
1 \arrow[r] & \mathfrak{z} \arrow[r, "{z\mapsto(z,z^{-1})}"'] & \widetilde{Z}\times\widetilde{Z} \arrow[r] & T_G \arrow[r] & 1
\end{tikzcd}

\marginal{NB on a supposé $c$ surjectif (cas où les « conditions d'isomorphie » sont
vérifiées)}

\struck{déjà le diagramme n'est commutatif,} \struck{mais ce diagramme} \ill{} \ill{}
l'inclusion $Z\subset T_G$ \ill{} \ill{} (connue sur le sous-groupe $c(\widetilde{Z})$ de $Z$)
comme induite par passage au quotient à partir de l'inclusion diagonale
$\widetilde{Z} \hookrightarrow \widetilde{Z}\times\widetilde{Z}$. Mais pour que cette
inclusion \struck{\ill{}} \add{envoie bien} $\mathfrak{z} = \operatorname{Ker} c$ dans
$\mathfrak{z} = \operatorname{Ker}(\widetilde{Z}\times\widetilde{Z}\to T_G)$, il faut (et il
suffit) qu'on ait
\[
(131)\qquad z^2 = 1 \qquad \forall z\in\mathfrak{z} ,
\]
et dans le premier carré de (130) est bien commutatif (l'hom $\mathfrak{z}\to\widetilde{Z} \times\widetilde{Z}$ de la deuxième ligne \struck{\ill{}} s'écrit aussi, \struck{\ill{}}
$z\mapsto(z,z)$ i.e. est \struck{\ill{}} le composé
\[
\mathfrak{z} \xrightarrow{\;\text{diag}\;} \mathfrak{z}\times\mathfrak{z}
\xrightarrow[\text{incl. can}]{} \widetilde{Z}\times\widetilde{Z}\ ).
\]
La commutativité du deuxième carré de (130) s'écrit sous la forme
\[
(132)\qquad c(\widetilde{\lambda}) = u_0(\widetilde{\lambda})\,u_1(\widetilde{\lambda})
\qquad \forall\widetilde{\lambda}\in\widetilde{Z}
\]
(qui pour $\widetilde{\lambda} = z\in\mathfrak{z}$ se réduit à (131)). On voit immédiatement que
les deux membres de (132) \struck{\ill{}} ont \uncertain{même} effet sur $(s_0,t_0)$ :
\[
(s_0,t_0) \longmapsto (\lambda s_0, \lambda t_0), \qquad \text{où}\quad \lambda =
c(\widetilde{\lambda})\,),
\]

\page{260}
\note{p. 133 de l'auteur.}

donc ils sont congrus dans $T_G$ modulo $\mathfrak{z}$. Pour voir qu'ils sont égaux, il
\struck{suffit de} suffit que leur image dans $G/H\simeq\widetilde{Z}$ soient les mêmes,
puisque l'on a $\mathfrak{z}\cap H = \{1\}$, i.e. que leurs images par $\widetilde{\delta}$
soient les mêmes, ce qui s'écrit sous la forme
\[
(133)\qquad c'(c(\widetilde{\lambda})) = \widetilde{\lambda}^2 \qquad
\forall\widetilde{\lambda}\in\widetilde{Z}
\]
(à rapprocher de (127), qui implique que les deux membres de (133) sont égaux mod
$\mathfrak{z}$).
\note{« (127) » : lecture incertaine du renvoi, peut-être (126).}

En résumé
\[
(134)\qquad \forall\widetilde{\lambda}\in\widetilde{Z}, \text{ on a}\qquad
\underset{Z\subset T_G}{\underset{\cap}{c(\widetilde{\lambda})}} \equiv
u_0(\widetilde{\lambda})\,u_1(\widetilde{\lambda}) \mod \mathfrak{z} .
\]
Pour que cette congruence soit une égalité, il faut et il suffit que l'on ait (133), (qui
implique (131)).

Quand de plus \struck{\ill{}} $c : \widetilde{Z}\to Z$ est surjectif, la formule (133) redonne
une façon d'expliciter le plongement de $Z$ dans $T_G$, en termes de la seule donnée du
diagramme (126).

En résumé, pour reconstituer tout le diagramme (64) : \uncertain{nous} pris — en même temps
\uncertain{que} (113), — à partir des données (89)
\note{la phrase se poursuit au-delà de ce lot ; le renvoi « (89) » est écrit par-dessus un
autre numéro.}

\end{document}
